% ============================================================
%  SECCIÓN — PRESUPUESTO (MURO)
% ============================================================
\section{Presupuesto de ejecución}
\label{sec:muro_presupuesto}

% -----------------------------------------------------------
\subsection{Mediciones}
\label{subsec:mediciones}
% -----------------------------------------------------------

Las mediciones se obtienen directamente de la geometría adoptada
(sección~\ref{sec:muro_predim}) y de los cálculos de materiales
del proceso constructivo (sección~\ref{sec:muro_proceso}):

\begin{align*}
  V_{\mathrm{exc}}    &= (B + 2 \times 0{,}60) \times (h_z + e_{\mathrm{HM}} + D_f)
                          \times L
                       = 6{,}80 \times 1{,}90 \times 150
                       = \mathbf{1\,938~\text{m}^3} \\
  V_{\mathrm{HM-15}}  &= B \times e_{\mathrm{HM}} \times L
                       = 5{,}60 \times 0{,}10 \times 150
                       = \mathbf{84~\text{m}^3} \\
  V_{\mathrm{HA-30}}  &= \underbrace{5{,}60 \times 0{,}80}_{V_{\mathrm{zapata}}}
                         \times 150
                       + \underbrace{\tfrac{0{,}60+0{,}75}{2} \times 6{,}50}
                                    _{V_{\mathrm{fuste}}} \times 150
                       = 672 + 659 = \mathbf{1\,331~\text{m}^3} \\
  P_{\mathrm{acero}}  &= 207~\text{kg/ml} \times 150~\text{ml}
                       = \mathbf{31\,050~\text{kg}} \\
  S_{\mathrm{enc,fuste}} &= 2 \times 6{,}50 \times 150
                          = \mathbf{1\,950~\text{m}^2} \\
  S_{\mathrm{enc,zap}}   &= 2 \times 0{,}80 \times 150
                          = \mathbf{240~\text{m}^2} \\
  S_{\mathrm{imp}}    &= 6{,}50 \times 150 = \mathbf{975~\text{m}^2} \\
  V_{\mathrm{rell}}   &= B_{\mathrm{talón}} \times H_f \times L
                       = 4{,}05 \times 6{,}50 \times 150
                       = \mathbf{3\,949~\text{m}^3} \\
  L_{\mathrm{pil}}    &= N_\text{pil} \times L_\text{pil}
                       = 150 \times 8{,}0
                       = \mathbf{1\,200~\text{m}}
                       \quad \text{(1 pilote/ml, $D=0{,}45$~m, hasta arcillas margosas)}
\end{align*}

% -----------------------------------------------------------
\subsection{Precios unitarios de referencia}
\label{subsec:precios_unitarios}
% -----------------------------------------------------------

\begin{table}[H]
  \centering
  \caption{Cuadro de precios unitarios de referencia (BEDEC / PREOC 2024)}
  \label{tab:precios}
  \begin{tabular}{@{}lcc@{}}
    \toprule
    \textbf{Concepto} & \textbf{P.U.} & \textbf{Ud.} \\
    \midrule
    Excavación a máquina en terreno cohesivo   & 3{,}00   & \euro/m³ \\
    Hormigón de limpieza HM-15, e=10 cm        & 80{,}00  & \euro/m³ \\
    Hormigón HA-30 puesto en obra (inc.\ bomba)& 105{,}00 & \euro/m³ \\
    Acero corrugado B~500~SD (inc.\ ferrallado)& 1{,}40   & \euro/kg \\
    Encofrado metálico doble cara (inc.\ grúa) & 65{,}00  & \euro/m² \\
    Impermeabilización HDPE $e \geq 2$~mm      & 18{,}00  & \euro/m² \\
    Lámina drenante nodular + geotextil        & 8{,}00   & \euro/m² \\
    Tubo dren flexible Ø~100~mm + grava        & 8{,}00   & \euro/ml \\
    Relleno y compactación trasdós (zahorra)   & 2{,}50   & \euro/m³ \\
    Pilotes CPI $D=0{,}45$~m (perf.\ + arm.\ + horm.) & 85{,}00 & \euro/m \\
    \bottomrule
  \end{tabular}
\end{table}

% -----------------------------------------------------------
\subsection{Cuadro de precios descompuesto}
\label{subsec:presupuesto}
% -----------------------------------------------------------

\begin{table}[H]
  \centering
  \caption{Presupuesto de ejecución del muro de contención ($L = 150$~ml)}
  \label{tab:presupuesto_muro}
  \resizebox{\textwidth}{!}{%
  \begin{tabular}{@{}clcccc@{}}
    \toprule
    \textbf{N.º} & \textbf{Descripción} & \textbf{Ud.}
    & \textbf{Medición} & \textbf{P.U. (\euro)} & \textbf{Total (\euro)} \\
    \midrule
    1 & Excavación a máquina                & m³ & 1\,938  &   3{,}00  &    5\,814 \\
    2 & Hormigón de limpieza HM-15          & m³ &    84   &  80{,}00  &    6\,720 \\
    3 & Hormigón HA-30 (zapata + fuste)     & m³ & 1\,331  & 105{,}00  &  139\,755 \\
    4 & Acero B~500~SD en armaduras         & kg & 31\,050 &   1{,}40  &   43\,470 \\
    5 & Encofrado fuste (doble cara)        & m² & 1\,950  &  65{,}00  &  126\,750 \\
    6 & Encofrado zapata (caras laterales)  & m² &   240   &  65{,}00  &   15\,600 \\
    7 & Impermeabilización HDPE trasdós     & m² &   975   &  18{,}00  &   17\,550 \\
    8 & Lámina drenante + geotextil         & m² &   975   &   8{,}00  &    7\,800 \\
    9 & Tubo dren Ø~100~mm                  & ml &   150   &   8{,}00  &    1\,200 \\
   10 & Relleno y compact.\ trasdós         & m³ & 3\,949  &   2{,}50  &    9\,873 \\
   11 & Pilotes CPI $D=0{,}45$~m, $L=8$~m  & m  & 1\,200  &  85{,}00  &  102\,000 \\
    \midrule
    \multicolumn{5}{r}{\textbf{PRESUPUESTO DE EJECUCIÓN MATERIAL (PEM)}}
      & \textbf{476\,532} \\
    \multicolumn{5}{r}{Gastos generales (13\,\%)}
      & 61\,949 \\
    \multicolumn{5}{r}{Beneficio industrial (6\,\%)}
      & 28\,592 \\
    \midrule
    \multicolumn{5}{r}{\textbf{Subtotal sin IVA}}
      & \textbf{567\,073} \\
    \multicolumn{5}{r}{IVA (21\,\%)}
      & 119\,085 \\
    \midrule
    \multicolumn{5}{r}{\textbf{PRESUPUESTO DE EJECUCIÓN POR CONTRATA (PEC)}}
      & \textbf{686\,158} \\
    \bottomrule
  \end{tabular}}
\end{table}

\begin{notabox}[Coste unitario por ml de muro]
  PEM por metro lineal de muro (incluyendo pilotes):
  $476\,532 \div 150 = \mathbf{3\,177~\text{\euro/ml}}$.
  Las partidas más significativas son el \textbf{hormigón HA-30} (29\,\% del PEM),
  el \textbf{encofrado del fuste} (27\,\%), los \textbf{pilotes CPI} (21\,\%) y el
  \textbf{acero} (9\,\%).
  La partida de pilotaje supone el \textbf{21\,\%} del coste total, justificada
  por la necesidad de reducir el asiento de $s = 7{,}51$~cm a valores admisibles
  (sección~\ref{subsec:asientos_muro}). Aun así, el coste total del muro con pilotes
  es significativamente inferior al de la solución de muro pantalla
  (sección~\ref{sec:muro_opciones}).
\end{notabox}
